\documentclass[10pt,a4paper]{amsart}
\usepackage[margin=19mm]{geometry}
\usepackage{amsmath,amssymb,mathtools}
\usepackage[scr=rsfs,cal=euler]{mathalfa}
\usepackage{xfrac}
\usepackage[unicode,hidelinks,bookmarks=false]{hyperref}
\newcommand{\ie}{i.e.\ }
\newcommand{\cf}{cf.\ }
\newcommand{\resp}{respectively\ }
\newcommand{\Subsection}{\subsection}
\providecommand{\nobreakdash}{\nobreak\hbox{-}}
\setlength{\emergencystretch}{2em}
\multlinegap0pt
\usepackage{dsfont}
\usepackage{amssymb}
\usepackage{caption}
\usepackage{tikz}
\newtheorem{theorem}{Theorem}
\newtheorem{definition}{Definition}
\newtheorem{lemma}{Lemma}
\newcommand{\eps}{\varepsilon}
\newcommand{\kint}{\vint}
\def\vint_#1{\mathchoice%
          {\mathop{\kern 0.2em\vrule width 0.6em height 0.69678ex depth -0.58065ex
                  \kern -0.8em \intop}\nolimits_{\kern -0.4em#1}}%
          {\mathop{\kern 0.1em\vrule width 0.5em height 0.69678ex depth -0.60387ex
                  \kern -0.6em \intop}\nolimits_{#1}}%
          {\mathop{\kern 0.1em\vrule width 0.5em height 0.69678ex
              depth -0.60387ex
                  \kern -0.6em \intop}\nolimits_{#1}}%
          {\mathop{\kern 0.1em\vrule width 0.5em height 0.69678ex depth -0.60387ex
                  \kern -0.6em \intop}\nolimits_{#1}}}
\title{Tug-of-war games related to $p$-Laplace type equations with zeroth order terms}
\author{Jeongmin Han}
\date{Source: arXiv:2506.07537v3 (CC BY 4.0)}
\begin{document}
\maketitle

\noindent\emph{Source and licence.} Tug-of-war games related to $p$-Laplace type equations with zeroth order terms. Version arXiv:2506.07537v3 (CC BY 4.0), distributed by arXiv under the Creative Commons Attribution 4.0 International licence, http://creativecommons.org/licenses/by/4.0/, as recorded in the arXiv licence metadata for this exact version and shown on the abstract page https://arxiv.org/abs/2506.07537v3. Full licence notice: \texttt{licenses/arxiv-2506.07537-CC-BY-4.0.txt}.\par\smallskip

\noindent\textit{Editorial scope.} Counted unit: the article's theorem conv (source0.tex lines 877--884). The article's complete original proof of it is delivered with it (source0.tex lines 885--998, one \texttt{proof} environment of 114 lines). No mathematics is written, repaired or re-derived: the statement and the proof are the article's own lines, in the article's own notation, and no retained line is substituted or re-flowed.  Retained uncounted as the source-local context the proof needs: lines 201--204; lines 208--213; lines 408--414; lines 535--539; lines 746--748; lines 771--778; lines 840--860; lines 865--873.  Nothing was omitted from any retained span, so the package carries no omission record.  No retained line contains a citation, so the wrapper carries no bibliography and no bibliography entry.  No retained line contains a figure environment, an image-inclusion command or drawing code, so no image is delivered and the package declares no asset dependency.  Editorial formatting only, in addition: an AMS \texttt{amsart} preamble that restates the article's own declarations and macros used by the retained lines, namely the environments \texttt{theorem}, \texttt{definition}, \texttt{lemma} on separate counters, together with the standard packages those lines need.  The retained environments print in this document as Theorem 1, Definition 1, Lemma 1 in order of appearance, whereas the article prints the same items with the numbering of its own full document.  The complete native source of the article as distributed in the arXiv e-print for this exact version is bundled unchanged as \texttt{sources/179538/source0.tex}.
\begin{align}\label{dpp_damp}
u_{\eps}(x)=(1-\gamma \eps^2)\bigg\{ \frac{\alpha}{2}\bigg(\sup_{B_\eps(x)}u_{\eps}+\inf_{B_\eps(x)}u_{\eps} \bigg)
+\beta\kint_{B_\eps(x)}u_{\eps}(y)dy \bigg\}
\end{align}

\begin{align}\label{eqmiin}
\left\{ \begin{array}{ll}
\Delta_{p}^{N} u-(p+n)\gamma u=0 & \textrm{in $ \Omega$,}\\
 u = F & \textrm{on $ \partial \Omega$}\\
\end{array} \right.
\end{align}

\begin{align}\label{deft} \mathcal{T}_{\eps}u(x) =
\left\{ \begin{array}{ll}
(1-\gamma \eps^2)\bigg\{ \frac{\alpha}{2}\big(\sup_{B_\eps(x)}u +\inf_{B_\eps(x)}u  \big)
+\beta\kint_{B_\eps(x)}u (y)dy \bigg\}& \textrm{if $ x\in \Omega$,}\\
 u(x)& \textrm{if $ x\in \Gamma_{\eps}$.}\\
\end{array} \right.
\end{align}

\begin{theorem}\label{intes}Assume that $B_{2}(0) \subset \Omega$. 
Suppose that $u_{\eps}$ satisfies \eqref{dpp_damp} for $ 0< \alpha < 1 $ and $\eps>0$ with $\gamma\eps^2<\frac{1}{2}$.
Then we have
$$ |u_{\eps} (x) - u_{\eps} (z) | \le C ||u_{\eps}||_{L^{\infty}(B_2)} ( |x-z| + \eps) ,$$ where $x, z \in B_{1}(0)$ and $C>0$ only depends on $ \alpha, \gamma$ and $n$.
\end{theorem}

\begin{align} \label{bdlip}
|F(x)-F(z)| \le L_F|x-z|
\end{align}

\begin{theorem}\label{bdes}
Assume that $\Omega$ satisfies the exterior sphere condition and $F$ satisfies \eqref{bdlip}.
Then for the value function $u_{\eps}$ with boundary data $F$, we have
\begin{align}\label{nearbd} 
& |u_{\eps}(x)-u_{\eps}(z)|  \le C(n,  R/\delta, \gamma)||F||_{C^{0,1}(\Gamma_{\eps})}( |x_0-y|+o(1)) +||F||_{C^{0,1}(\Gamma_{\eps})}\delta
\end{align} 
for any $x\in\Omega$ and $z\in \Gamma_{\eps}$.
\end{theorem}

\begin{definition} \label{vissol} 
A function $u \in C(\Omega) $ is a viscosity solution to \eqref{eqmiin} if
the following conditions hold:
\begin{itemize}
\item[(a)] for all $ \varphi \in C^{2}(\Omega ) $ touching $u$ from above at $x_{0} \in \Omega $, 
\begin{align*} 
\left\{ \begin{array}{ll}
 \Delta_{p}^{N}\varphi(x_{0})  \ge (p+n)\gamma \varphi(x_{0}) \qquad \qquad \textrm{if $ D\varphi (x_{0}) \neq 0 $,}\\
\lambda_{\max}((p-2)D^{2}\varphi(x_{0}) )  + \Delta\varphi(x_{0}) \ge (p+n)\gamma \varphi(x_{0})  \qquad \textrm{if $D\varphi (x_{0}) = 0   $.}\\
\end{array} \right. 
\end{align*}
\item[(b)] for all $ \varphi \in C^{2}(\Omega ) $ touching $u$ from below at $x_{0} \in \Omega $, 
\begin{align*} 
\left\{ \begin{array}{ll}
 \Delta_{p}^{N}\varphi(x_{0})  \le (p+n)\gamma \varphi(x_{0}) \qquad \qquad \textrm{if $ D\varphi (x_{0}) \neq 0 $,}\\
\lambda_{\min}((p-2)D^{2}\varphi(x_{0}) )  + \Delta\varphi(x_{0}) \le (p+n)\gamma \varphi(x_{0})  \qquad \textrm{if $D\varphi (x_{0}) = 0   $.}\\
\end{array} \right. 
\end{align*}
\end{itemize}
Here, the notations $\lambda_{\max}(X)$ and $\lambda_{\min}(X)$ mean the largest and the smallest eigenvalues of a symmetric matrix $X$. 
\end{definition}

\begin{lemma} \label{aras}
Let $ \{ u_{\eps} : \overline{\Omega} \to \mathbb{R}, \eps>0 \} $ be a set of functions such that
\begin{itemize}
\item[(a)] there exists a constant $C >0$ so that $|u_{\eps}(x)| < C $ for every $ \eps > 0$ and every $x \in \overline{\Omega} $. 
\\ \item[(b)] given $ \eta > 0$, there are constants $r_{0}$ and $\eps_{0}$ so that for every $\eps >0$ and $x,z \in \overline{\Omega}$ with $|x-z| < r_{0} $, it holds
$$ |u_{\eps}(x)-u_{\eps}(z)| < \eta .$$
\end{itemize}
Then, there exists a uniformly continuous function $u: \overline{\Omega} \to \mathbb{R} $ and a subsequence $\{ u_{\eps_{\textrm{I}}} \} $ such that $ u_{\eps_{\textrm{I}}}  $ uniformly converges to $u$  in $\overline{\Omega}$, as $i \to \infty$.
\end{lemma}

\begin{theorem}\label{conv}Assume that $\Omega$ satisfies the exterior sphere condition and  $ F \in C^1(\Gamma_{\eps}) $
satisfies \eqref{bdlip}.
Let $u_{\eps}$ denote the solution to \eqref{dpp_damp} with boundary data $ F  $ for each $\eps>0$.
Then, there exist a function $ u: \overline{\Omega}_{\eps} \to \mathbb{R}$ and a subsequence 
$ \{ \eps_{\textrm{I}} \} $ such that 
$$ u_{\eps_{\textrm{I}}} \to u \qquad \textrm{uniformly in} \quad \overline{\Omega}.$$
Furthermore, $u$ is a viscosity solution to \eqref{eqmiin}.
\end{theorem}

\begin{proof}
We first observe the existence of a uniform convergent subsequence of $\{u_{\eps}\}$.
By using the definition of $u_{\eps} $, we have
 $$||u_{\eps} ||_{L^{\infty}(\Omega)} \le ||F||_{L^{\infty}(\Omega)} < \infty $$
for any $\eps > 0$ and thus $u_{\eps} $ are uniformly bounded. 
We also see the equicontinuity in the sense of Lemma \ref{aras} by Theorem \ref{intes} and Theorem \ref{bdes},
and hence we can apply the Arzel\`{a}-Ascoli criterion.
Therefore, we can find a uniformly convergent subsequence, still denoted by $u_{\eps}$, to a function $u \in C(\overline{\Omega})$.

Next, we verify that $u$ is a viscosity solution to \eqref{eqmiin}.
We first observe that
$$ u(x) = \lim_{\eps \to 0}u_{\eps}(x) = F(x)$$
for any $x \in \partial\Omega$.
Thus, it is sufficient to show that $u$ is a solution to 
$$\Delta_{p}^{N} u-(p+n)\gamma u=0   \qquad \textrm{in} \ \Omega$$
in the viscosity sense.
Without loss of generality, it is enough to prove that $u$ satisfies (b) in Definition \ref{vissol}. 

We fix $x \in \Omega$ and consider a small neighborhood of $x$ such that 
$ B_{R}(x)  \subset \Omega $ for some $R>0$.
For each $\eps>0$, we can consider a function $\varphi \in C^{2}(B_{R}(x))$ touching $u$ from below at $x$.
Since $u_{\eps}$ converges uniformly to $u$, for sufficiently small $\eps >0$, we can find a convergent sequence $\{x_{\eps}\}$
satisfying the following property 
$$ (u_{\eps}-\varphi)(z) \ge (u_{\eps}-\varphi)(x_{\eps})-\eta_{\eps} \qquad\textrm{for any} \ \ z\in B_{R}(x) $$
for some $x_{\eps}\in B_{R}(x)$ and small $\eta_{\eps}>0$. 
We observe that $x_{\eps} \to x$ as $ \eps \to 0$ in that case.

Recall \eqref{deft}. For $x \in \Omega$,
we have
\begin{align*}u_{\eps}(x)=
\mathcal{T}_{\eps}u_{\eps}(x) =(1-\gamma \eps^2)\bigg\{ \frac{\alpha}{2}\bigg(\sup_{B_\eps(x)}u_{\eps} +\inf_{B_\eps(x)}u_{\eps}  \bigg)
+\beta\kint_{B_\eps(x)}u_{\eps} (y)dy \bigg\}.
\end{align*}
Set $ \psi = \varphi +( u_{\eps}-\varphi)(x_{\eps})$.
Then we observe that
$u_{\eps} \ge \psi - \eta_{\eps}$ in $B_{R}(x) $.
By using this, it follows that
\begin{align*}
u_{\eps}(x_{\eps})=\mathcal{T}_{\eps}u_{\eps}(x_{\eps}) 
 \ge \mathcal{T}_{\eps}\big(\psi(x_{\eps})-\eta_{\eps}) =\mathcal{T}_{\eps}\psi(x_{\eps})-
 (1-\gamma \eps^2)\eta_{\eps}
\end{align*}
and 
\begin{align*}
 \mathcal{T}_{\eps}\psi(x_{\eps}) =\mathcal{T}_{\eps}\varphi(x_{\eps})+(1-\gamma \eps^2)( u_{\eps}-\varphi)(x_{\eps}).
\end{align*}
Hence, we get  
\begin{align*}
u_{\eps}(x_{\eps}) \ge \mathcal{T}_{\eps}\varphi(x_{\eps})+(1-\gamma \eps^2)\big(( u_{\eps}-\varphi)(x_{\eps})-\eta_{\eps}\big)
\end{align*}
and this yields 
\begin{align} \label{2est1} \gamma\eps^2  (u_{\eps}-\varphi)(x_{\eps})+(1-\gamma \eps^2)\eta_{\eps} \ge \mathcal{T}_{\eps}\varphi (x_{\eps})-\varphi (x_{\eps}). 
\end{align}
%We remark that the first inequality holds since $\varphi$ touches $u_{\eps}$ from below.

By the Taylor expansion, we can compute that 
\begin{align*} (1-\gamma \eps^2)&\bigg\{ \frac{\alpha}{2}\bigg(\sup_{B_\eps(x)}\varphi  +\inf_{B_\eps(x)}\varphi   \bigg)
+\beta\kint_{B_\eps(x)}\varphi  (y)dy \bigg\}
\\&= \varphi(x) + \frac{1}{2}\bigg( \frac{\Delta_{p}^{N} \varphi(x)}{p+n}-\gamma \varphi(x) \bigg)\eps^2+o(\eps^2).
\end{align*}
Since we have assumed that $\varphi$ has the $C^2$-regularity, $\varphi$ attains its local minimum at a point $  x_{\eps}^1 $ in $\overline{B}_\eps(x)$, that is,
$$\varphi(x_{\eps}^1)=\inf_{B_{\eps}(x)}\varphi .$$
Assume that $x_{\eps}^1 \neq x$. Then for $\tilde{x}_{\eps}^1=2x-x_{\eps}^1$ (the mirror point of $x_{\eps}^1$ with respect to $x$), we have 
\begin{align*}
& (1-\gamma \eps^2)\bigg\{ \frac{\alpha}{2}\bigg(\sup_{B_\eps(x)}\varphi  +\inf_{B_\eps(x)}\varphi   \bigg)
+\beta\kint_{B_\eps(x)}\varphi  (y)dy \bigg\} - \varphi (x) \\
 & \ge  (1-\gamma \eps^2)\bigg\{ \frac{\alpha}{2}\big(\varphi(x_{\eps}^1)  +  \varphi(\tilde{x}_{\eps}^1)  \big)
+\beta\kint_{B_\eps(x)}\varphi  (y)dy \bigg\} - \varphi (x) 
\\ & = \frac{1}{2} \eps^{2} \bigg(  \frac{(p-2) \langle D^{2} \varphi (x) \nu_{\eps}, \nu_{\eps} \rangle+\Delta\varphi (x)  }{p+n} - \gamma  \varphi (x)\bigg)+o(\eps^2),
\end{align*}
where $\nu_{\eps}=(x-x_{\eps}^1)/|x-x_{\eps}^1|$. 
By \eqref{2est1}, we obtain 
\begin{align} \begin{split} \label{2est2}
 \gamma\eps^2  (u_{\eps}-\varphi)(x_{\eps})&+(1-\gamma \eps^2)\eta_{\eps} \\&\ge   
 \frac{1}{2} \eps^{2} \bigg(  \frac{(p-2) \langle D^{2} \varphi (x) \nu_{\eps}, \nu_{\eps} \rangle+\Delta\varphi (x)  }{p+n} - \gamma \varphi (x)\bigg)+o(\eps^2)  .
\end{split}
\end{align}

Suppose that $D\varphi (x) \neq 0 $. Since $x_{\eps} \to x$ when $\eps \to 0$, it follows that
$$\nu_{\eps}  \to -\frac{D\varphi (x)}{|D\varphi (x)|} $$
as $\eps \to 0$ (we remark that we can find a proper subsequence of $\{\nu_{\eps}\}$ in that case due to $D\varphi (x) \neq 0 $).
and this implies 
$$(p-2) \langle D^{2} \varphi (x) \nu_{\eps}, \nu_{\eps} \rangle+\Delta\varphi (x) \to \Delta_{p}^{N} \varphi (x)   .$$
Meanwhile, we also have
$$(u_{\eps}-\varphi)(x_{\eps}) \to (u-\varphi)(x) =0$$
by the uniform convergence of $u_{\eps}$. 
We take $\eta=o(\eps^2)$ and
divide both side in \eqref{2est2} by $\eps^{2}$.
By letting $\eps \to 0$, we finally get
$$ 0 \ge \Delta_{p}^{N} \varphi (x) -(p+n)\gamma \varphi (x)  .$$



On the other hand, if $D\varphi (x) = 0  $,
we observe that 
\begin{align*} 
& (1-\gamma \eps^2)\bigg\{ \frac{\alpha}{2}\bigg(\sup_{B_\eps(x)}\varphi  +\inf_{B_\eps(x)}\varphi   \bigg)
+\beta\kint_{B_\eps(x)}\varphi  (y)dy \bigg\} - \varphi (x) \\
 & \ge  (1-\gamma \eps^2)\bigg\{ \frac{\alpha}{2}\big(\varphi(x_{\eps}^1)  +  \varphi(\tilde{x}_{\eps}^1)  \big)
+\beta\kint_{B_\eps(x)}\varphi  (y)dy \bigg\} - \varphi (x) 
\\ & \ge \frac{1}{2} \eps^{2} \bigg(  \frac{(p-2) \lambda_{\min}(D^{2}\varphi(x))+\Delta\varphi (x)  }{p+n} - \gamma  \varphi (x)\bigg)+o(\eps^2).
\end{align*}
By using $x_{\eps} \to x$ as $\eps \to 0$ and the continuity of the map $z  \mapsto \lambda_{\min}(D^{2}\varphi(z)) $, we can verify that
\begin{align} \label{2est3} 0 \ge \frac{1}{p+n} \big\{ \Delta\varphi(x) + (p-2)\lambda_{\min}(D^{2}\varphi(x) ) \big\}
-\gamma \varphi (x)
\end{align}
by a similar computation in the previous case.



We can also show the inequality in the opposite direction by considering a function $\varphi $ touching $u$ from above
and using a similar argument.
Then we complete the proof.
\end{proof}

\end{document}
